\subsection{Item-Level Analysis}

\subsubsection{Cross-Item Patterns and Interpretation Framework}

We examine several patterns that emerge across the 14 PSTRE items from the set of selected actions. These cross-item regularities --- how the directional balance and sparsity of the selected actions relate to the functional structure of each task --- both characterize the selection results and motivate a unified framework for interpreting the action effects for each item. We first describe two such patterns and then set out the interpretive framework they imply, which is applied throughout the real data analysis.

\paragraph*{Pattern 1: Directional balance varies with task structure.}
The directional balance between positive and negative selected actions varies across items in a manner that reflects the functional structure of the task. Several items show a predominantly positive pattern --- for example, PS1-4 (Sprained Ankle Part~1; 6 positive, 1 negative) and PS2-3 (Book Order; 2 positive, 0 negative) --- indicating that correct problem-solving in those tasks is characterized by a distinct set of productive behaviors. Other items show a negative-dominant or balanced pattern --- such as PS2-5 (Reply All; 3 positive, 7 negative) and PS2-4 (Meeting Room; 6 positive, 8 negative) --- indicating that behaviors more prevalent among incorrect respondents also carry discriminating information.

\paragraph*{Pattern 2: Selection sparsity reflects task structure, not vocabulary size.}
The number of selected actions does not increase monotonically with the size of the action vocabulary. PS2-2 (Club Membership Part~2) has the largest vocabulary (289 unique actions) yet yields only 4 selected actions, whereas PS1-1 and PS1-6 have smaller vocabularies (119 and 182, respectively), each producing 16 selected actions. This pattern indicates that the sparsity of behaviorally important actions depends on the task structure and the degree to which the interface generates discriminating interactions, rather than on the vocabulary size alone.

\paragraph*{Interpretation framework.}
The patterns above motivate a unified framework for interpreting the action effects for each item. The key quantity is $\Delta E_{jl}$, the difference between the group-level action effects for correct and incorrect respondents defined in Equations~\eqref{eq:group_effect}--\eqref{eq:delta_E}. Two interpretive levels are distinguished throughout.

At the \emph{first-order} level, the selection of an action as important, that is, having a 95\% HPD interval of $\Delta E_{jl}$ that excludes zero, indicates that its latent representation carries information about response accuracy beyond what is explained by respondent ability $\alpha_i$ and item difficulty $\beta_j$. This inference is available for every selected action regardless of sequence richness.

At the \emph{second-order} level, the sign and magnitude of $\Delta E_{jl}$ admit a process-level interpretation. Because $\Delta E_{jl}$ is a conditional partial effect --- the residual association between the action's latent representation and response accuracy after the conditions of $\alpha_i$ and $\beta_j$ --- it does not measure the total effect of performing the action. Specifically, the marginal frequency advantage that a given action may exhibit for correct respondents is largely absorbed by the ability parameter $\alpha_i$ during estimation. The residual effect $\Delta E_{jl}$ therefore reflects \emph{within-ability-group} differences in the sequential context in which the action was performed, as encoded by the LSTM autoencoder. This absorption mechanism produces a systematic discordance: actions performed predominantly by correct respondents can carry negative $\Delta E_{jl}$ and actions performed more often by incorrect respondents can carry positive $\Delta E_{jl}$. Such discordance should not be interpreted as paradoxical; rather, it is a direct consequence of the ability-adjusted nature of the partial effect.

Second-order interpretation is applicable only when the action sequences are sufficiently rich to provide latent representations that encode process-level context beyond the identity of individual actions. In items with extended, multi-step sequences, the LSTM encoder captures contextual information --- such as exploration order, trial-and-error patterns, and temporal pacing --- that is partly independent of $\alpha_i$, and the sign of $\Delta E_{jl}$ supports substantive behavioral inference. In items with short, simple sequences where the only meaningful behavioral variation is the identity of the final choice, $\bar{C}_{ijl}$ co-varies strongly with $\alpha_i$, and second-order interpretation is unreliable.

Based on the functional character of their selected actions, the 14 PSTRE items are classified into three types. In \textit{outcome-dominant} items, most selected actions correspond to final-choice operations; interpretation is accordingly restricted to the first-order level. In \textit{process-dominant} items, selected actions are distributed across intermediate behavioral steps, supporting both first- and second-order interpretation throughout the solution process. In \textit{hybrid} items, both outcome-proximal and process-level actions appear among the selected set, permitting first-order inference for all selected actions and second-order inference for the intermediate ones. Table~\ref{tab:item_types} assigns each of the 14 PSTRE items to one of the three types and records the defining functional characteristic of each.

\begin{table}[htb]
\centering
{\small
\begin{tabular}{cll}
\toprule
Label & Description & Items \\
\midrule
Outcome-dominant & Most selected actions are final-choice operations & PS1-3, PS1-4, PS1-5, PS1-7, PS2-3 \\
Process-dominant & Selected actions span intermediate process steps & PS1-1, PS1-2, PS2-1, PS2-5, PS2-6 \\
Hybrid           & Outcome-proximal and process actions coexist     & PS1-6, PS2-2, PS2-4, PS2-7 \\
\bottomrule
\end{tabular}
}
\caption{Typology of the 14 PSTRE items by the functional character of their selected important actions, giving the defining feature and the member items of each of the three types.}
\label{tab:item_types}
\end{table}

The following subsections present one representative case study per type: PS1-3 (CD Tally) for outcome-dominant items, PS1-1 (Party Invitations Part~1) for process-dominant items, and PS2-7 (Lamp Return) for hybrid items. Detailed results for the remaining items are provided in Section~D of the Supplementary Material.

\subsubsection{Outcome-Dominant Items: CD Tally (PS1-3)}

The CD Tally task asks respondents to count the number of Blues-genre compact discs listed in a sortable spreadsheet and enter the total using a numeric dropdown menu. Meaningful behavioral variation is concentrated almost entirely at the final selection step, with limited discriminating information in intermediate navigation actions. Three actions were identified as important: the correct answer entry, a sorting strategy, and an interface switch. At the first-order level, all three carry information about response accuracy beyond what is explained by respondent ability and item difficulty. 

\begin{table}[htb]
\centering
\small
\begin{tabular}{clrrl}
\toprule
Dir. & Action & $n_C$ & $n_I$ & $\Delta E$ (95\% HPD) \\
\midrule
$(+)$ & \texttt{combobox\_menulist\_index9}     & 490 &   2 & $\phantom{-}$0.2726 (0.2457, 0.3009) \\
$(+)$ & \texttt{combobox\_sortablecol1\_index3} & 199 &  20 & $\phantom{-}$0.0005 (0.0000, 0.0010) \\
$(-)$ & \texttt{toolbar\_webapp}               & 479 & 196 & $-$0.0037 ($-$0.0064, $-$0.0008) \\
\bottomrule
\end{tabular}

\caption{Important actions for PS1-3 (CD Tally). $n_C$ and $n_I$ denote the number of respondents in the correct and incorrect groups who performed the action at least once.}
\label{tab:ps13_actions}
\end{table}

The action \texttt{combobox\_menulist\_index9} corresponds to selecting 9 as the Blues CD count, which is the correct answer; it was performed by 490 correct respondents and only 2 incorrect respondents. This extreme frequency imbalance constitutes quasi-complete separation in the logistic component and the resulting $\Delta E = 0.273$ is likely inflated in magnitude. The direction of the effect is interpretable but the magnitude should not be compared directly with those of other actions.

The action \texttt{combobox\_sortablecol1\_index3} corresponds to sorting the spreadsheet by the genre column, a strategy that facilitates counting by grouping all Blues titles together. Among the three selected actions, this is the only one that admits a partial second-order interpretation: it represents a deliberate procedural choice rather than an outcome-defining operation and its positive $\Delta E = 0.0005$ suggests that this sorting strategy is associated with a modestly favorable sequential context after controlling for respondent ability.

The action \texttt{toolbar\_webapp} represents switching from the spreadsheet interface to the web application, an operation not required to complete the task. Although it was performed more frequently by correct respondents ($n_C = 479$) than by incorrect respondents ($n_I = 196$), its $\Delta E = -0.004$ is negative. As described in the interpretation framework above, this discordance arises because the marginal frequency advantage is absorbed by $\alpha_i$; the residual partial effect indicates that, within ability-matched groups, unnecessary interface switching is associated with a less favorable sequential context.

This item illustrates both the strengths and the interpretive boundaries of the proposed framework for outcome-dominant tasks. The model identifies the correct answer selection without requiring explicit labeling of which interface element corresponds to the correct response and it detects a strategic sorting behavior that contributes process-level information. At the same time, the predominance of final-choice actions and the quasi-complete separation in the correct answer action limit the extent to which the partial effects support detailed process-level inference.


% Supp의 해석으로 넘기기
% Another outcome-dominant item, PS1-5 (Sprained Ankle Part~2), exhibits a more extreme version of this limitation. All four selected actions in PS1-5 are combobox choices corresponding to specific answer options, and no intermediate process action was identified as important. In this item, the correct answer action (\texttt{combobox\_index2}; $n_C = 661$, $n_I = 10$) carries a negative $\Delta E = -0.035$ despite its overwhelming frequency advantage for correct respondents. This systematic discordance between marginal frequency and partial effect across all selected actions arises because the simple procedural structure of the task---respondents view web pages and then select an answer---limits the latent representation to the identity of the chosen response, which co-varies strongly with respondent ability. For items with this structure, interpretation is accordingly restricted to the first-order level: the selected actions carry information beyond ability and difficulty, but the sign of $\Delta E$ does not support stable inference about behavioral strategies.

\subsubsection{Process-Dominant Items: Party Invitations (PS1-1)}

The Party Invitations task requires respondents to read incoming email messages and sort each sender into either a "Can Come" or "Cannot Come" folder based on the stated availability. The task is scored polytomously (0--3), with the binary outcome defined as a score of 2 or 3. Sixteen actions were identified as important tied for the largest count among all 14 items. The action sequences for this item are sufficiently rich (median sequence length 16; IQR 11--25) to support both first- and second-order interpretation throughout.

\begin{table}[htb]
\centering
{\small
\begin{tabular}{clrrl}
\toprule
Dir. & Action & $n_C$ & $n_I$ & $\Delta E$ (95\% HPD) \\
\midrule
$(+)$ & \texttt{mail-drop\_u01a-cancomefolder}     & 783 & 130 & $\phantom{-}$0.3263 ($\phantom{-}$0.2121, $\phantom{-}$0.4439) \\
$(+)$ & \texttt{mail-drag\_u01a-item104}           & 772 &  58 & $\phantom{-}$0.1372 ($\phantom{-}$0.0759, $\phantom{-}$0.1956) \\
$(+)$ & \texttt{mail-drop\_u01a-cannotcomefolder}  & 662 &  33 & $\phantom{-}$0.1029 ($\phantom{-}$0.0494, $\phantom{-}$0.1554) \\
$(+)$ & \texttt{mail-drop\_u01a-partyfolder}       &  58 &  70 & $\phantom{-}$0.1044 ($\phantom{-}$0.0351, $\phantom{-}$0.1771) \\
$(+)$ & \texttt{mail-drag\_u01a-item201}           &  34 &  88 & $\phantom{-}$0.0273 ($\phantom{-}$0.0104, $\phantom{-}$0.0445) \\
$(+)$ & \texttt{folder-viewed\_u01a-cannotcomefolder} & 346 & 188 & $\phantom{-}$0.0000 ($\phantom{-}$0.0000, $\phantom{-}$0.0000) \\
$(-)$ & \texttt{folder-viewed\_u01a-inboxfolder}  & 334 & 150 & $-$0.0232 ($-$0.0351, $-$0.0110) \\
$(-)$ & \texttt{toolbar\_mailapp}                 &  17 &  54 & $-$0.0004 ($-$0.0007, $-$0.0001) \\
\bottomrule
\end{tabular}
}
\caption{Selected important actions for PS1-1 (Party Invitations Part~1). Eight representative actions are shown; the complete list of all 16 selected actions is provided in Section~D of the Supplementary Material. $n_C$ and $n_I$ as in Table~\ref{tab:ps13_actions}.}
\label{tab:ps11_actions}
\end{table}

At the first-order level, the 16 selected actions encompass the core operations required for the task --- reading target messages (e.g., \texttt{mail-viewed\_item105}, \texttt{mail-viewed\_item102}), dragging messages for classification (e.g., \texttt{mail-drag\_item104}, \texttt{mail-drag\_item101}), and dropping messages into the designated folders (\texttt{mail-drop\_cancomefolder}, \texttt{mail-drop\_cannotcomefolder}) --- as well as non-task-relevant exploration behaviors such as viewing extraneous messages (\texttt{mail-viewed\_item202}), opening the parent folder (\texttt{folder-viewed\_partyfolder}), and switching interfaces (\texttt{toolbar\_mailapp}). The selection of both task-relevant and non-task-relevant actions provides evidence that the model captures process-level information extending beyond the identity of the correct classification operations. The action \texttt{folder-viewed\_u01a-cannotcomefolder} is excluded from the second-order discussion below, as its estimated effect is negligible ($\Delta E = 0.000$; 95\% HPD: 0.000, 0.000).

The largest positive effects are associated with the core classification steps: dropping messages into the Can Come folder ($\Delta E = 0.326$), dragging a task-relevant message ($\Delta E = 0.137$), and dropping messages into the Cannot Come folder ($\Delta E = 0.103$). For these actions, the marginal frequency pattern and the partial effect are consistent, indicating that the sequential context in which these actions were performed provides incremental information beyond the contribution already absorbed by respondent ability.

Several selected actions exhibit the discordance pattern described in the interpretation framework. The actions \texttt{mail-drag\_u01a-item201} ($n_C = 34$, $n_I = 88$; $\Delta E = +0.027$) and \texttt{mail-drag\_u01a-item202} ($n_C = 36$, $n_I = 84$; $\Delta E = +0.002$) involve messages that are not classification targets and both were performed more often by incorrect respondents. Their positive residual partial effects may reflect that, among respondents of comparable ability, those who engaged with extraneous messages as part of a systematic exploration pattern experienced a more favorable sequential context than those who did so in a less structured manner. This interpretation is consistent with the sequential context encoded by the LSTM but should be regarded as tentative given the relatively small number of correct respondents who performed these actions.

The action \texttt{folder-viewed\_u01a-inboxfolder} ($n_C = 334$, $n_I = 150$; $\Delta E = -0.023$) illustrates the reverse discordance: performed more often by correct respondents, yet carrying a negative partial effect. Opening the inbox folder is the default starting point for the task and is performed by a large proportion of all respondents, so its positive marginal frequency association with correct response is absorbed by $\alpha_i$. The negative residual partial effect suggests that, within ability-matched groups, repeated returns to the inbox may reflect less efficient navigation behavior---re-checking messages already read rather than proceeding with classification.

The action \texttt{mail-drop\_u01a-partyfolder} ($n_C = 58$, $n_I = 70$; $\Delta E = +0.104$) represents placing a message into the parent Party folder rather than the required subfolders. Its positive partial effect suggests that, among respondents of comparable ability, those who used the parent folder did so in a sequential context associated with eventual task success. For example, as a temporary holding step before redistributing messages to the appropriate subfolders. This interpretation is plausible but not directly verifiable from the model output alone.

Overall, the extended action sequences in this item allow the LSTM encoder to capture contextual information that is partially independent of respondent ability and the resulting partial effects reveal behavioral distinctions --- such as the difference between systematic and unsystematic exploration of non-target messages --- that would not be visible from action frequencies or response outcomes alone.

\subsubsection{Hybrid Items: Lamp Return (PS2-7)}

The Lamp Return task requires respondents to initiate a product return by navigating between a simulated email client and a web form: the correct procedure involves requesting a return authorization number via an email link, retrieving the number from the subsequent reply, and submitting the completed return request through the web interface. Eight actions were identified as important spanning both intermediate procedural steps and final confirmatory operations. The action sequences are moderately long (median sequence length 15; IQR 6--24), providing sufficient sequential context for both first- and second-order interpretation.

\begin{table}[htb]
\centering
\small
\begin{tabular}{clrrl}
\toprule
Dir. & Action & $n_C$ & $n_I$ & $\Delta E$ (95\% HPD) \\
\midrule
$(-)$ & \texttt{textlink\_u023-pg2\_txt21\_self}    & 252 &  13 & $-$0.0032 ($-$0.0046, $-$0.0018) \\
$(-)$ & \texttt{combobox\_u023-pg2\_menu1\_index1}  & 541 &  18 & $-$0.0171 ($-$0.0329, $-$0.0005) \\
$(-)$ & \texttt{textbox-onfocus\_u023-pg2\_txt19}  & 519 &   9 & $-$0.0340 ($-$0.0472, $-$0.0224) \\
$(+)$ & \texttt{mail-viewed\_u23-item305}          & 480 & 132 & $\phantom{-}$0.0010 ($\phantom{-}$0.0003, $\phantom{-}$0.0016) \\
$(+)$ & \texttt{toolbar\_back-btn}                & 502 & 233 & $\phantom{-}$0.0353 ($\phantom{-}$0.0166, $\phantom{-}$0.0545) \\
$(+)$ & \texttt{button\_u023-pg2\_txt22}           & 512 &  11 & $\phantom{-}$0.0505 ($\phantom{-}$0.0397, $\phantom{-}$0.0615) \\
$(+)$ & \texttt{keypress2}                        & 345 &   9 & $\phantom{-}$0.4619 ($\phantom{-}$0.1941, $\phantom{-}$0.7227) \\
$(-)$ & \texttt{textlink\_u023-default\_txt5}      & 150 & 324 & $-$0.0030 ($-$0.0055, $-$0.0005) \\
\bottomrule
\end{tabular}
\caption{Important actions for PS2-7 (Lamp Return). $n_C$ and $n_I$ as in Table~\ref{tab:ps13_actions}. Actions are listed in approximate procedural order.}
\label{tab:ps27_actions}
\end{table}

At the first-order level, the eight selected actions cover each major stage of the multi-step return procedure: requesting the authorization email (\texttt{textlink\_txt21\_self}), selecting the return reason (\texttt{combobox\_menu1\_index1}), activating the input field for the authorization number (\texttt{textbox-onfocus\_txt19}), reviewing the authorization email (\texttt{mail-viewed\_item305}), navigating back to the return form (\texttt{toolbar\_back-btn}), entering the authorization number (\texttt{keypress2}), and submitting the final request (\texttt{button\_txt22}). The breadth of this selection indicates that the model identifies discriminating information at multiple points along the procedural sequence rather than concentrating it at a single outcome-defining step.

At the second-order level, a systematic sign pattern emerges when the selected actions are arranged in approximate procedural order. The three actions associated with the earlier stages of the procedure --- requesting the authorization email ($n_C = 252$, $n_I = 13$), selecting the return reason ($n_C = 541$, $n_I = 18$), and activating the authorization number input field ($n_C = 519$, $n_I = 9$) --- all carry negative $\Delta E$ despite being performed predominantly by correct respondents. This is an instance of the discordance pattern described in the interpretation framework: the strong positive marginal frequency associations at these early stages are absorbed by $\alpha_i$, and the residual partial effects reflect within-ability-group differences in the sequential context surrounding these actions.

By contrast, the four actions associated with the later stages --- reviewing the authorization email ($\Delta E = +0.001$), navigating back to the return form ($\Delta E = +0.035$), entering the authorization number ($\Delta E = +0.462$), and submitting the final request ($\Delta E = +0.051$) --- all carry positive partial effects. This anterior--posterior sign differentiation can be interpreted as follows. The early procedural steps are initiated by a large proportion of respondents, including many who attempt the return procedure but do not complete it; among ability-matched respondents, reaching these early stages does not by itself discriminate those who ultimately succeed. The later stages, by contrast, are reached primarily by respondents who sustain the multi-step procedure to completion, and the positive partial effects at these stages suggest that completing the later procedural steps carries an incremental association with task success not fully explained by respondent ability. This pattern is consistent with the interpretation that correct and incorrect respondents diverge not at the point of initiating the procedure but at the point of completing it.

Two individual actions warrant additional comment. The action \texttt{keypress2}, corresponding to typing the return authorization number, carries the largest $\Delta E$ among the eight selected actions ($\Delta E = 0.462$; 95\% HPD: 0.194, 0.723); however, the extreme frequency imbalance ($n_C = 345$, $n_I = 9$) constitutes quasi-complete separation, and the wide HPD interval reflects the resulting estimation uncertainty. The direction of the effect is interpretable, but the magnitude should not be compared with those of other actions. The action \texttt{textlink\_u023-default\_txt5} ($n_C = 150$, $n_I = 324$; $\Delta E = -0.003$) was performed predominantly by incorrect respondents and carries a negative partial effect; because its specific interface function could not be conclusively identified from the available log metadata, we refrain from offering a substantive behavioral interpretation.

This item illustrates the analytical value of the hybrid classification. The coexistence of intermediate process actions and final confirmatory operations within the same selected set allows the model to locate where in a multi-step procedure the behavioral trajectories of correct and incorrect respondents diverge, and the systematic anterior--posterior sign pattern provides a more informative characterization of the problem-solving process than would be available from either final-step analysis or action frequency comparison alone.